%==============================================================
% Taliqa Muhib — Curriculum Vitae
% Tailored for: Data Scientist, Care ADHD
% Compile with: pdflatex Taliqa_Muhib_CV_CareADHD.tex
%   (or upload to Overleaf and press Recompile)
%==============================================================
\documentclass[a4paper,10pt]{article}

\usepackage[T1]{fontenc}
\usepackage[utf8]{inputenc}
\usepackage[margin=1.5cm]{geometry}
\usepackage{enumitem}
\usepackage{titlesec}
\usepackage{xcolor}
\usepackage{hyperref}
\usepackage{fontawesome5}
\usepackage{multicol}

% ---- Colours
\definecolor{accent}{HTML}{0F8F86}
\definecolor{ink}{HTML}{10233B}
\definecolor{muted}{HTML}{4A5B70}

\hypersetup{
  colorlinks=true,
  urlcolor=accent,
  linkcolor=accent
}

% ---- Section styling
\titleformat{\section}
  {\color{accent}\scshape\large\bfseries}
  {}{0em}{}[\vspace{-0.6em}\color{accent}\rule{\linewidth}{0.8pt}]
\titlespacing{\section}{0pt}{8pt}{6pt}

\setlist[itemize]{leftmargin=1.2em, itemsep=1pt, topsep=2pt}

\newcommand{\entry}[4]{%
  \noindent\textbf{#1}\hfill{\color{muted}\small #2}\\[-0.2em]
  {\color{accent}\small #3}\\[0.1em]
  #4
}

\pagestyle{empty}
\setlength{\parindent}{0pt}

%==============================================================
\begin{document}

% ---------------- Header ----------------
\begin{center}
  {\Huge\bfseries\color{ink} Taliqa Muhib}\\[0.35em]
  {\color{accent}\bfseries HealthTech Data Scientist $\cdot$ Patient Pathway Analytics $\cdot$ Predictive Modelling}\\[0.5em]
  {\small
    \faEnvelope\ \href{mailto:taliqa.muhib@gmail.com}{taliqa.muhib@gmail.com} \quad
    \faPhone\ +44 7732 477820 \quad
    \faMapMarker*\ Pakistan-based --- open to UK relocation / hybrid
  }\\[0.3em]
  {\small
    \faGlobe\ \href{https://talikamuhib.github.io/Talikamuhib/}{Portfolio} \quad
    \faGithub\ \href{https://github.com/Talikamuhib}{github.com/Talikamuhib} \quad
    \faLinkedin\ \href{https://www.linkedin.com/in/talika-muhib}{linkedin.com/in/talika-muhib} \quad
    \faChartBar\ \href{https://public.tableau.com/app/profile/taliqa.muhib1820}{Tableau Public}
  }
\end{center}

% ---------------- Summary ----------------
\section*{Professional Summary}
HealthTech Data Scientist specialising in \textbf{patient pathway analytics, predictive
modelling, and clinical service performance}. I turn healthcare, operational, and product
data into actionable insights that improve patient outcomes and service efficiency. MRes
research background, strong Python / SQL / Tableau toolkit, and a responsible, GDPR-aware
approach to sensitive clinical data. Passionate about improving access to ADHD and
neurodiversity care through data.

% ---------------- Skills ----------------
\section*{Core Skills}
\begin{itemize}
  \item \textbf{Data Science \& ML:} Python, Pandas, NumPy, Scikit-learn, Statistical Analysis,
        Predictive Modelling, Classification, Regression, Feature Engineering, Model Evaluation
  \item \textbf{Healthcare Analytics:} Patient Pathway Analysis, Appointment Attendance
        Prediction, Referral Trends, Waiting-Time Analysis, Clinical Data, Healthcare
        Operations, GDPR-aware Analysis
  \item \textbf{Data \& BI:} SQL, SQL Server, MySQL, Tableau, Power BI, Excel, Dashboards,
        KPI Reporting, Data Quality Checks
  \item \textbf{Research \& AI:} EEG Analysis, Graph Signal Processing, Medical Imaging,
        Vision Transformers, Healthcare AI, Responsible AI, LLM Evaluation
\end{itemize}

% ---------------- Projects ----------------
\section*{Selected Projects}

\entry{ADHD Patient Pathway Analytics Dashboard}
      {\href{https://public.tableau.com/views/MonthlyReferralADHDanalysis/Dashboard1}{Live dashboard}}
      {Python $\cdot$ SQL $\cdot$ Pandas $\cdot$ Tableau $\cdot$ Healthcare Analytics}
      {\begin{itemize}
        \item Built an end-to-end analytics dashboard modelling \textbf{1,200 referrals} to
              analyse demand, waiting times, assessment progression, and appointment attendance
              across an ADHD care pathway.
        \item Surfaced a \textbf{17\% did-not-attend (DNA) rate} and \textbf{77\%
              18-week-target compliance}, identifying DNA reduction as a lever for recovering
              wasted clinical capacity.
        \item Showed urgent cases correctly prioritised ($\sim$30 days vs $\sim$85 for routine);
              combines Python/SQL pipelines with Tableau to help service leads understand demand
              and bottlenecks --- designed to support service planning, not clinical diagnosis.
      \end{itemize}}
\vspace{0.5em}

\entry{ADHD Appointment No-Show (DNA) Prediction}
      {\href{https://github.com/Talikamuhib}{Code / Streamlit app}}
      {Python $\cdot$ Scikit-learn $\cdot$ Streamlit $\cdot$ Classification}
      {\begin{itemize}
        \item Developed an interactive risk model (Logistic Regression, \textbf{ROC-AUC $\approx$ 0.62})
              predicting which referrals are most likely to no-show, as decision-support for
              targeted reminders and scheduling optimisation.
        \item Handled class imbalance and avoided target leakage; built a Streamlit app with
              live prediction, ROC curve, feature importance, and confusion-matrix views.
      \end{itemize}}
\vspace{0.5em}

\entry{EEG Neurodegenerative Disease Detection (MRes Research)}
      {Coventry University}
      {Python $\cdot$ EEG $\cdot$ Graph Signal Processing $\cdot$ Machine Learning}
      {\begin{itemize}
        \item Used EEG connectivity, graph signal processing, and ML to support early-detection
              research for neurodegenerative disease.
        \item Preprocessed EEG signals, built graph-based connectivity features, and trained
              reproducible ML pipelines --- research-only, not intended for diagnosis.
      \end{itemize}}
\vspace{0.5em}

\entry{Kidney Stone Detection using Vision Transformers}
      {Final-year project}
      {Python $\cdot$ Vision Transformers $\cdot$ Medical Imaging $\cdot$ Deep Learning}
      {\begin{itemize}
        \item Applied Vision Transformer models to CT scan images for exploratory kidney-stone
              detection, studying model performance on medical imaging as a research prototype.
      \end{itemize}}
\vspace{0.5em}

\entry{Health LLM Evaluation Framework}
      {Research}
      {LLM Evaluation $\cdot$ Healthcare AI $\cdot$ Responsible AI $\cdot$ Benchmarking}
      {\begin{itemize}
        \item Built a responsible-AI framework evaluating healthcare language-model responses
              across accuracy, safety, bias, and clinical reliability, with defined benchmarks
              and scoring --- not intended for clinical use.
      \end{itemize}}

% ---------------- Education ----------------
\section*{Education}
\textbf{MRes, Computer Science} --- Coventry University\\
{\color{muted}\small Research: early detection of neurodegenerative disease using EEG, graph signal processing, and machine learning.}\\[0.4em]
\textbf{BS, Computer Science}\\
{\color{muted}\small Final-year project: kidney-stone detection from CT images using Vision Transformers.}

% ---------------- Other work ----------------
\section*{Other Work}
\begin{itemize}
  \item \textbf{EdTech Learning Platform} --- coding and digital-skills training to widen access to technical education.
  \item \textbf{Glacier \& Climate Monitoring} --- data-driven monitoring of glacial change for environmental analysis.
  \item \textbf{Portfolio Analytics Dashboard} --- general KPI reporting and data-visualisation practice (Tableau, Power BI).
  \item \textbf{Interactive Game Projects} --- small front-end builds exploring logic, state, and interactivity.
\end{itemize}

% ---------------- Strengths ----------------
\section*{Key Strengths}
\begin{itemize}
  \item Translate complex analysis into clear dashboards and recommendations for non-technical clinical and operational stakeholders.
  \item Responsible, GDPR-aware handling of sensitive data; treats models as decision support, not clinical judgement.
  \item Commercially minded --- frames analytics around service improvement, capacity, and patient outcomes.
\end{itemize}

\vfill
{\footnotesize\color{muted}\itshape Portfolio projects use synthetic demonstration data and are decision-support prototypes, not intended for clinical diagnosis or deployment.}

\end{document}
